\scanpage{54}
% @ line 1 / text
L'\deleted{application} morphisme)
% @ lines 2–3 / display
\[
 \begin{aligned}
 D&\longmapsto D\cap X=Y\\
 \check P=\operatorname{Dr}(P)&\xrightarrow{\sim}\operatorname{Div}_2^+(X)
 \end{aligned}\tag*{(22)}
\]
% @ lines 4–6 / text
est un iso, qui associe à \deleted{\unsure{tt}} $H\subset\widetilde{\unsure{\mathfrak g|g}}$ le diviseur $Y\subset X$ des \deleted{droites}\inserted{ss-fibrés inv. $\underline L$} dans $\widetilde{\mathfrak g}$ qui \uline{à la fois} sont dans $H$, et satisfont $q|\underline L=0$.
% @ lines 7–10 / text
\deleted{C'est donc} Comme schéma sur $S$, il représente donc le sous-foncteur de $\unsure{\mathbb P^{*}|\check{\mathbb P}}(H)=\check{\mathbb P}(H)$ défini par l'annulation de $q$. On trouve donc, grâce à la \deleted{\unsure{proposition}}\inserted{prop. précédente}, un iso. can.
% @ line 11 / display
\[
 Y\simeq\operatorname{Spec} A\tag*{(23)}
\]

% @ lines 12–16 / text
Dans l'interprétation $X\xleftarrow{\sim}\mathbb P(V)$, via $L\rightarrow\underline L$, on voit que $Y$ est le sous-schéma de $\mathbb P(V)$ formé des $L$ tels que $AL\subset L$ ou encore (si $A=\mathcal O_S1+\mathcal O_Su$) tel que $uL\subset L$, ou encore, si $u$ est inv. i.e. section de $\check{\mathbb V}(M^*)$
% @ below line 16, left / deletion
\deleted{\illegible{1}}
% @ lines 17–18 / display
\[
 Y=X^u\qquad\begin{aligned}
 &\text{sous-schéma des}\\[-2pt]
 &\text{invariants sous }u
 \end{aligned}\tag*{(24)}
\]
% @ lines 19–24, continues on next page / text
On retrouve réciproquement $A$ quand on connaît $Y$, en notant que \deleted{l'opération de}\inserted{l'isomorphisme} (23) est fonctoriel, donc compatible aux autom provenant de $u\in\deleted{\unsure{A}}\Gamma(M^*)$ (i.e. ici $\{u:V\to V)$, en faisant opérer $u$ par autom int. sur $M^*$, et de la façon évidente sur $X$ \unsure{$\leftarrow$|\ldots}
% archive sheet number: 53
